\uuid{siNp}
\titre{Étude de fonction}
\auteur{bodin ; hanani ; mzari}
\organisation{exo7}
\licence{CC BY-NC-SA 4.0}
\datecreate{2020-01-20}
\niveau{L1}
\module{Analyse}
\chapitre{Fonctions circulaires et hyperboliques inverses}
\sousChapitre{Autre}
\difficulte{2}
\theme{Fonctions usuelles, étude de fonction}
\exo7qcmid{lille-1.fonctions-usuelles.020}

\contenu{
\begin{question}
\qclassification{126.00}
\qtags{motcle=étude de fonction}
Soit $f(x)=x+ \sqrt{ 1-x^2}$.  Quelles sont les assertions vraies ?

\begin{answers}
    \good{Le domaine de définition de $f$ est $[-1,1]$.}

    \bad{$f$ est croissante sur $[-1,1]$.}

    \bad{$f$ établit une bijection de $[0,1]$ dans $[1,\sqrt 2]$.}

    \good{$f$ établit une bijection de $[-1,\frac{1}{\sqrt 2}]$ dans $[-1,\sqrt 2]$.}
\end{answers}
\begin{explanations}
$f$ est définie sur $[-1,1]$, dérivable sur $]-1,1[$ et $f'(x)= 1-\frac{x}{\sqrt{1-x^2}}$. En étudiant les variations de $f$, on déduit que $f$ établit une bijection de $[-1,\frac{1}{\sqrt 2}]$ dans $[-1,\sqrt 2]$ et de $[\frac{1}{\sqrt 2}, 1]$  dans $[1,\sqrt 2]$.
\end{explanations}
\end{question}
}
